% ===== PRÉAMBULE CANONIQUE — à coller VERBATIM en tête de CHAQUE .tex =====
\documentclass[11pt,a4paper]{article}
\usepackage[utf8]{inputenc}\usepackage[T1]{fontenc}\usepackage{lmodern}
\usepackage[french]{babel}
\usepackage{amsmath,amssymb,mathtools}\usepackage{icomma}
\usepackage{siunitx}\sisetup{locale=FR}  % \qty{}{}, \si{}, \ang{}
\usepackage{xcolor}\usepackage{colortbl}
\usepackage{tikz}\usepackage{pgfplots}\pgfplotsset{compat=1.18}
\usepackage[siunitx,europeanresistors]{circuitikz} % circuits (élec.)
\usepackage[most]{tcolorbox}\usepackage{enumitem}\usepackage{array,booktabs}
\usepackage[a4paper,margin=2.2cm]{geometry}
\usetikzlibrary{arrows.meta,calc,angles,quotes,positioning,patterns,%
  backgrounds,shapes.geometric,decorations.pathmorphing,decorations.markings,intersections}
\definecolor{primary}{RGB}{222,49,99}\definecolor{primarydark}{RGB}{150,22,55}
\definecolor{defcol}{RGB}{0,114,178}\definecolor{methcol}{RGB}{52,92,148}
\definecolor{excol}{RGB}{0,135,90}\definecolor{attncol}{RGB}{200,40,55}
\tcbset{boxbase/.style={enhanced,breakable,boxrule=0.9pt,arc=2.5pt,
  left=3mm,right=3mm,top=2.3mm,bottom=2.3mm,fonttitle=\bfseries,coltitle=white}}
% --- boîtes TUTORIEL (noms EXACTS, ne pas renommer) ---
\newtcolorbox{cle}[1][]{boxbase,colback=defcol!6,colframe=defcol,
  title={\textsc{À retenir}\ifx\relax#1\relax\relax\else\textmd{ : \itshape #1}\fi}}
\newtcolorbox{methode}[1][]{boxbase,colback=methcol!7,colframe=methcol,sharp corners,
  title={\textsc{Méthode}\ifx\relax#1\relax\relax\else\textmd{ --- \itshape #1}\fi}}
\newtcolorbox{exemplebox}[1][]{boxbase,colback=excol!5,colframe=excol,
  title={\textsc{Exemple}\ifx\relax#1\relax\relax\else\textmd{ : \itshape #1}\fi}}
\newtcolorbox{piege}[1][]{boxbase,colback=attncol!6,colframe=attncol,
  title={\textsc{Piège}\ifx\relax#1\relax\relax\else\textmd{ : \itshape #1}\fi}}
\newtcolorbox{remarque}[1][]{boxbase,colback=black!4,colframe=black!45,
  title={\textsc{Remarque}\ifx\relax#1\relax\relax\else\textmd{ : \itshape #1}\fi}}
% --- boîtes EXERCICE / CORRIGÉ (noms EXACTS, auto-comptées) ---
\newtcolorbox[auto counter]{exo}[1][]{boxbase,colback=primary!4,colframe=primary,
  title={\textsc{Exercice}~\thetcbcounter\ifx\relax#1\relax\relax\else\textmd{ --- \itshape #1}\fi}}
\newtcolorbox[auto counter]{corrige}[1][]{boxbase,colback=excol!4,colframe=excol,
  title={\textsc{Corrigé}~\thetcbcounter\ifx\relax#1\relax\relax\else\textmd{ --- \itshape #1}\fi}}
\newcommand{\vv}[1]{\vec{#1}}\newcommand{\ds}{\displaystyle}
\newcommand{\R}{\mathbb{R}}\newcommand{\N}{\mathbb{N}}\newcommand{\Z}{\mathbb{Z}}
% =========================================================================
\begin{document}
\sisetup{per-mode=symbol}

\begin{exo}[deuxième loi : $a=F/m$]
Un palet de masse $m=\qty{3}{\kilogram}$ glisse sans frottement sur une patinoire. On lui applique
une force résultante horizontale $\vv F$ d'intensité \qty{12}{\newton}.
\begin{center}
\begin{tikzpicture}[>=Stealth]
  \draw[thick,black!55] (-2.6,0) -- (2.8,0);
  \fill[primarydark] (-1.2,0) rectangle (-0.4,0.7);
  \node[white,font=\small] at (-0.8,0.35){$m$};
  \draw[->,line width=1.1pt,defcol] (-0.4,0.35) -- ++(1.7,0) node[right]{$\vv F=\qty{12}{\newton}$};
\end{tikzpicture}
\end{center}
\begin{enumerate}[label=\alph*)]
  \item Calcule l'accélération $a$ du palet.
  \item Si l'on double la force ($\qty{24}{\newton}$), que devient $a$ ?
\end{enumerate}
\end{exo}

\begin{exo}[deuxième loi : $F=ma$]
Un mobile de masse $m=\qty{2}{\kilogram}$ prend une accélération $a=\qty{5}{\metre\per\second\squared}$.
\begin{enumerate}[label=\alph*)]
  \item Quelle est l'intensité de la force résultante appliquée ?
  \item Vérifie l'unité obtenue à partir de $\si{\newton}=\si{\kilogram\metre\per\second\squared}$.
\end{enumerate}
\end{exo}

\begin{exo}[retrouver la masse]
Une force résultante de \qty{50}{\newton} communique à un objet une accélération de
$\qty{2,5}{\metre\per\second\squared}$. Quelle est la masse de l'objet ?
\end{exo}

\begin{exo}[première loi : la résultante]
Un livre est posé, immobile, sur une table. Il subit son poids $\vv G$ (vers le bas) et la réaction
normale $\vv N$ de la table (vers le haut).
\begin{center}
\begin{tikzpicture}[>=Stealth]
  \fill[black!10] (-1.8,-0.15) rectangle (1.8,0);
  \draw[thick,black!55] (-1.8,0) -- (1.8,0);
  \fill[primarydark] (-0.45,0) rectangle (0.45,0.6);
  \draw[->,line width=1.1pt,attncol] (0,0) -- ++(0,-1.2) node[below]{$\vv G$};
  \draw[->,line width=1.1pt,excol] (0,0.6) -- ++(0,1.2) node[above]{$\vv N$};
\end{tikzpicture}
\end{center}
\begin{enumerate}[label=\alph*)]
  \item Que vaut la résultante des forces sur le livre ?
  \item Qu'en déduit-on pour son accélération, et pour son état de mouvement ?
\end{enumerate}
\end{exo}

\begin{exo}[direction de l'accélération]
Un chariot est tiré par une unique force résultante horizontale $\vv F$ dirigée vers la droite.
\begin{enumerate}[label=\alph*)]
  \item Dans quelle direction et quel sens pointe le vecteur accélération $\vv a$ ?
  \item Cette accélération pourrait-elle pointer vers le haut ? Justifie par la $2^{\text{e}}$ loi.
\end{enumerate}
\end{exo}
\end{document}
